\chapter{大事年表}\label{app:timeline}

本附录汇总本书涉及的关键事件，供快速检索。可信度标注沿用正文制度：
\conf{A} 一手权威来源，\conf{B} 机构或平台公开信息，\conf{C} 社群考据或二手评述。
同一事件存在多个来源口径时并列列出，不作合并。

\section{正典奶龙：IP 与产业线}\label{app:timeline:canon}

\begin{longtable}{@{}p{3.4cm}p{8.6cm}p{2.2cm}@{}}
  \toprule
  时间 & 事件 & 可信度 \\ \midrule
  2017 & 谢承恩与团队扎根深圳，成立第七印象文化传媒（深圳）有限公司 & \conf{A} \\ \midrule
  2019 & 开始构思奶龙形象，定位为“传递治愈、欢乐的积极情感” & \conf{A} \\ \midrule
  2020-01-22 & 奶龙 IP 在短视频平台正式上线 & \conf{B} \\ \midrule
  2020 & 抖音粉丝破千万，单条视频 3 天点播破亿（第一版教材口径，未获独立权威源印证） & \conf{C} \\ \midrule
  2020--2021 & 短视频、有声读物专辑陆续上线；有声读物 8 个月播放量破亿（第一版教材口径） & \conf{C} \\ \midrule
  2023 & 第一季番剧《奶龙搞怪大作战》播出 & \conf{B} \\ \midrule
  2024 & 第二季番剧《奶龙与小七之大战暴暴龙》播出（第一版教材口径） & \conf{C} \\ \midrule
  2024 & 第七印象参评深圳市新型业态企业，获 25 万元扶持奖励；开始参加深圳文博会 & \conf{A} \\ \midrule
  2025-01 & 第三季番剧《奶龙之奇妙冒险》上线，一个月内全网播放量破 1 亿 & \conf{B} \\ \midrule
  2025 底 & 授权合作品牌超 200 家；衍生 SKU 逾 2000 个；动画译制 15 种语言；微信表情包累计发送量突破 80 亿次 & \conf{A/B} \\ \midrule
  2026-01-15 & 《人民日报》报道“一只‘奶龙’背后的创新链”（全网粉丝超 3800 万、流量破千亿、销往 57 个国家和地区） & \conf{A} \\ \midrule
  2026-01-19 至 25 & “奶龙六周年”生日会：线下观影、五城地标大屏（武汉、青岛、天津、成都、苏州）、三城门店见面会 & \conf{B} \\ \midrule
  2026-01-22 & 奶龙六周年；以“歌手”身份发布首支单曲《奶龙的小确幸》 & \conf{B} \\ \midrule
  2026-01-24 & 深圳海岸影城六周年包场观影，约 200 位粉丝观看《奶龙搞怪大作战 2》精选内容 & \conf{B} \\ \midrule
  2026-01-25 & 亮相中央广播电视总台（CMG）《动漫盛典》，以开场节目之一身份出现 & \conf{B} \\ \midrule
  2026-01-28 & 新番《奶龙搞怪大作战 2》在腾讯、优酷、爱奇艺等平台上线，覆盖风行、聚好看、银河少儿等大屏终端；嘉佳卡通、金鹰卡通等电视台陆续播出 & \conf{B} \\ \midrule
  荣誉（未标注具体届次时间） & 亮相昂西国际动画节；获第 21 届中美电视节“年度最佳动画剧集”金天使奖；入围第 30 届亚洲电视大奖“最佳儿童节目” & \conf{B} \\ \bottomrule
\end{longtable}

\section{民间奶蛙：起源与传播线}\label{app:timeline:naiwa}

\begin{longtable}{@{}p{3.4cm}p{8.6cm}p{2.2cm}@{}}
  \toprule
  时间 & 事件 & 可信度 \\ \midrule
  2024-03-07 & 有网友上传将“奶蛙”形象与抖音博主“季洪行”视频进行 AI 融合的内容；“奶蛙”长期被误认为与此有关 & \conf{C} \\ \midrule
  2024-10-02 & 星野 AI 平台用户“纲手【百豪】”创建“奶龙”智能体（ID \fname{zNCvQDKfXy}），平台 AI 生成的头像出现形体畸变，被社群认定为最初版奶蛙形象；该智能体现已无法访问，评论区关闭 & \conf{C} \\ \midrule
  2024-12 前后 & 抖音用户对“欢迎大家恐龙世界”的声讨截图中发现奶蛙照片，并由此出现“星野 AI 生成”的推测（截图含 AI 生成水印） & \conf{C} \\ \midrule
  2026（上半年，具体日期不详） & 抖音视频生成模型 Seedance 推出，奶蛙二创视频生产门槛显著降低（二手评述口径） & \conf{C} \\ \midrule
  2026-03 & 奶龙官方呼吁停止丑化二创，关闭视频评论区并批量拉黑发布奶蛙的账号 & \conf{C} \\ \midrule
  2026-08-04 & B 站用户“雾蚀心忘症”在星野 App 中找到符合形象的奶蛙智能体并发布评论，起源考据取得关键突破 & \conf{C} \\ \midrule
  2026-08 & 以 B 站 UP 主“看仙碎”为代表的一批爱好者推进起源考据，发布考据视频（视频号 BV1ndMk65EwM，第一版教材口径） & \conf{C} \\ \midrule
  2026 全年 & 奶蛙成为中文互联网高传播度迷因；被二手评述称为“2026 年传播最广、影响最深、二创最活跃的 meme” & \conf{C} \\ \bottomrule
\end{longtable}

\section{高校与线下活动线}\label{app:timeline:campus}

\begin{longtable}{@{}p{3.4cm}p{8.6cm}p{2.2cm}@{}}
  \toprule
  时间 & 事件 & 可信度 \\ \midrule
  2026-06 & 哈尔滨工业大学 2026 年毕业典礼暨学位授予仪式出现第一个广为人知的奶蛙立牌 & \conf{C} \\ \midrule
  2026-06 & 上海交通大学发起“流动奶蛙”活动（最早的流动奶蛙至今仍在用） & \conf{C} \\ \midrule
  2026-06 & 中国海洋大学“奶 OUC”四奶蛙立牌出现（第一个广为人知的复数奶蛙立牌） & \conf{C} \\ \midrule
  2026-06-24 & 北京航空航天大学奶龙同好会成立 & \conf{C} \\ \midrule
  2026-08 初 & 北航同好会人数突破 100 人 & \conf{C} \\ \midrule
  2026-09-28 & 北京航空航天大学学院路校区绿园投放流动奶蛙立牌；旁边随后出现丰川祥子立牌；当晚被保安带走，存活约 9 小时 & \conf{C} \\ \midrule
  2026-10 初 & 北航同好会人数突破 200 人（26 级新生入学后） & \conf{C} \\ \midrule
  2026（具体日期不详） & 浙江大学某地招生宣传立牌出现浙大奶蛙形象 & \conf{C} \\ \bottomrule
\end{longtable}

\noindent\srcnote{年表中“具体日期不详”的条目，
来源为第一版教材或社群考据词条，未给出确切日期；
本书不推测日期。2020 年“抖音粉丝破千万”等
未获独立权威源印证的条目，引用时请保留 \conf{C} 标注。}
